\chapter{Background and Related Work}
\label{chap:background}

\section{Diabetic Retinopathy and Fundus Images}
\label{sec:dr-background}
\writingnote{Explain the original five severity labels and the binary detection
task. Use cited examples of fundus images only where they help the reader
understand the input and labels. Do not equate every DR-present image with
referable or severe disease.}

\section{CNNs and Transfer Learning}
\label{sec:transfer-learning}
\writingnote{Introduce only the CNN concepts needed later: feature extraction,
the classification head, pretraining, and fine-tuning. Starter method references
are ResNet~\cite{he2016resnet} and EfficientNet~\cite{tan2019efficientnet}.
Add verified references for other architectures actually discussed.}

\section{Explainability and Grad-CAM}
\label{sec:explainability-background}
\writingnote{Explain the purpose and basic mechanism of
Grad-CAM~\cite{selvaraju2017gradcam}. Separate visual attribution from lesion
segmentation, causal explanation, and clinical validation. Reserve details of
our chosen target layer and implementation for Chapter~\ref{chap:methods}.}

\section{Related Studies and Positioning}
\label{sec:related-work}
\writingnote{Compare relevant DR studies using the fields in
Table~\ref{tab:related-work}. Explain differences in datasets, labels, and
validation protocols before discussing scores. Identify the scope of this work
without claiming an unverified research gap or state-of-the-art result.}

\tableplaceholder{Comparison framework for related DR studies.}{tab:related-work}{%
Populate with verified studies. Columns: source and year; dataset; binary or
multiclass task; method; evaluation protocol; metrics; explainability; limitations.
Split a wide table rather than reducing the text to an unreadable size.}
